\usepackage[utf8]{inputenc}
\usepackage[T1]{fontenc}
\usepackage[frenchb]{babel}
\usepackage[french]{minitoc}
\usepackage{hyperref}
\usepackage{titletoc}
%\usepackage{lmodern} % Pour changer le pack de police
\usepackage{amsmath}
\usepackage{amsthm}
\usepackage{amsfonts}
\usepackage{amssymb}
\usepackage{makeidx}
\usepackage{graphicx}
\usepackage{color}
\usepackage{xcolor}
\usepackage{pgf,tikz}
\usepackage{setspace}
\usepackage{soul}
%\usepackage{ulem}
%\usepackage{lmodern}
\usepackage{apalike}
%\usepackage{fourier}
%\usepackage[scaled=2]{helvet}
\usepackage{courier}
%\usepackage{eurosym}
%\usepackage{bookman}

%\usepackage{charter}

%\usepackage{newcent}
%\usepackage{mathpazo}

\usepackage{mathptmx}
\usepackage{url}
\usepackage{verbatim}
\usepackage{moreverb}
\usepackage{listings}
\usepackage{fancyhdr}
\usepackage{wrapfig}
\usepackage{colortbl}
\usepackage{wasysym}
\usepackage{lastpage}
\usepackage{listings}
%\usepackage{glossaries}
%\makeglossaries
\usepackage[top=2cm,left=2cm,right=2cm,bottom=3cm]{geometry}
\hypersetup{colorlinks,%
    linkcolor=black,
    %citecolor=green,%vert1,%
    urlcolor=cyan,%
}
%\usepackage[linesnumbered, french]{algorithm2e}
%\usepackage{fp}
%%\usepackage{pstricks-add}
%\RestyleAlgo{algoruled}
%\SetKwInput{donne}{Données}
%\SetKwInput{res}{Résultat}
%\SetKw{ret}{Retourner}
\usepackage{algorithmic}
\usepackage{algorithm}
\renewcommand{\algorithmicrequire}{\textbf{Données:}}
\renewcommand{\algorithmicensure}{\textbf{Résultat:}}
\renewcommand{\algorithmicwhile}{\textbf{Tant que}}
\renewcommand{\algorithmicdo}{\textbf{Faire}}
\renewcommand{\algorithmicendwhile}{\textbf{Fin Tant que}}
\renewcommand{\algorithmicend}{\textbf{Fin}}
\renewcommand{\algorithmicif}{\textbf{Si}}
\renewcommand{\algorithmicendif}{\textbf{Fin du si}}
\renewcommand{\algorithmicelse}{\textbf{Sinon}}
\renewcommand{\algorithmicelsif}{\textbf{fin du sinon}}
\renewcommand{\algorithmicthen}{\textbf{alors}}
\renewcommand{\algorithmicfor}{\textbf{Pour}}
\renewcommand{\algorithmicforall}{\textbf{pour tout}}
\renewcommand{\algorithmicto}{\textbf{à}}
\renewcommand{\algorithmicendfor}{\textbf{Fin Pour}}
\renewcommand{\algorithmicdo}{\textbf{faire}}
\renewcommand{\algorithmicloop}{\textbf{boucler}}
\renewcommand{\algorithmicendloop}{\textbf{fin de la boucle}}
\renewcommand{\algorithmicrepeat}{\textbf{répéter}}
\renewcommand{\algorithmicuntil}{\textbf{jusqu’à}}
\renewcommand{\algorithmicprint}{\textbf{afficher}}
\newtheorem{theoreme}{Théorème}[section]
\newtheorem{theoremes}{Théorèmes}[section]
\newtheorem{TheoremeDefinition}{Théorème et Définition}[section]
\newtheorem{definition}{Définition}[section]
\newtheorem{definitions}{Définitions}[section]
\newtheorem{proposition}{Proposition}[section]
\newtheorem{propostions}{Propositions}[section]
\newtheorem{corollaire}{Corollaire}[section]
\newtheorem{corollaires}{Corollaires}[section]
\newtheorem{lemme}{Lemme}[section]
\newtheorem{lemmes}{Lemmes}[section]
\newtheorem{remarque}{Remarque}[section]
\newtheorem{exemple}{Exemple}[section]
\newtheorem{formulemoyenne}{Formule de la moyenne}
\newtheorem*{formulmoydiscr}{caption}
\newcommand{\R}{ \mathbb{R} }
\lstset{
	literate=
	{á}{{\'a}}1 {é}{{\'e}}1 {í}{{\'i}}1 {ó}{{\'o}}1 {ú}{{\'u}}1
	{Á}{{\'A}}1 {É}{{\'E}}1 {Í}{{\'I}}1 {Ó}{{\'O}}1 {Ú}{{\'U}}1
	{à}{{\`a}}1 {è}{{\`e}}1 {ì}{{\`i}}1 {ò}{{\`o}}1 {ù}{{\`u}}1
	{À}{{\`A}}1 {È}{{\'E}}1 {Ì}{{\`I}}1 {Ò}{{\`O}}1 {Ù}{{\`U}}1
	{ä}{{\"a}}1 {ë}{{\"e}}1 {ï}{{\"i}}1 {ö}{{\"o}}1 {ü}{{\"u}}1
	{Ä}{{\"A}}1 {Ë}{{\"E}}1 {Ï}{{\"I}}1 {Ö}{{\"O}}1 {Ü}{{\"U}}1
	{â}{{\^a}}1 {ê}{{\^e}}1 {î}{{\^i}}1 {ô}{{\^o}}1 {û}{{\^u}}1
	{Â}{{\^A}}1 {Ê}{{\^E}}1 {Î}{{\^I}}1 {Ô}{{\^O}}1 {Û}{{\^U}}1
	{œ}{{\oe}}1 {Œ}{{\OE}}1 {æ}{{\ae}}1 {Æ}{{\AE}}1 {ß}{{\ss}}1
	{ű}{{\H{u}}}1 {Ű}{{\H{U}}}1 {ő}{{\H{o}}}1 {Ő}{{\H{O}}}1
	{ç}{{\c c}}1 {Ç}{{\c C}}1 {ø}{{\o}}1 {å}{{\r a}}1 {Å}{{\r A}}1
	{€}{{\EUR}}1 {£}{{\pounds}}1
}
\lstdefinestyle{stylepython}{
	language=Python,
	basicstyle=\ttfamily,
	commentstyle=\color{green},
	keywordstyle=\color{violet},
	stringstyle=\color{olive},
	numberstyle=\tiny,
	numbers=left,
	stepnumber=1,
	numbersep=5pt,
	linewidth =18cm, 
	xleftmargin =2cm, 
	aboveskip =5mm, 
	belowskip =1cm,
    float =!h,
    frame =tBlR,
%    rulesep =1mm,
%    framesep =5mm,
    %framerule=2pt,
    xrightmargin=5mm,
    xleftmargin=5mm,
     rulecolor={\color[gray]{0.6}},rulesepcolor={\color[gray]{0.9}}
}
\newcommand{\N}{ \mathbb{N} }
\newcommand{\abC}{$\mathcal{C}^{0}([a,b])$}
